\documentclass[english,twoside, openright, hidelinks, 11 pt, class=report,crop=false]{standalone}
\input{../../preamb}
\input{../bm}

\begin{document}
\opgt

\nes \nes

\op{opgstatsentrodd1}
Gitt datasettet
\[ 2\quad12\quad 3\quad 0\quad 2\quad 5\quad 8\quad2\quad 11 \]
Finn \os
\abch{
	\item typetallet \item medianen \item gjennomsnittet 
}

\op{opgstatsentrodd2}
Gitt datasettet
	\[ 9\quad6\quad 3\quad 0\quad 8\quad 5\quad 8\quad4\quad 10\quad 8 \quad 6 \]
Finn \os
\abch{
\item typetallet \item medianen \item gjennomsnittet 
}

\op{opgstatsentrpar1}
Gitt datasettet
\[ 11\quad7\quad 16\quad 0\quad 8\quad 9\quad 8\quad5\quad 16\quad 5 \]
Finn \os
\abch{
	\item typetallet \item medianen \item gjennomsnittet 
}

\op{opgstatsentrpar2}
Gitt datasettet
\[ 6\quad11\quad 14\quad 5\quad 6\quad 9\quad 8\quad5\quad 11\quad 5\quad 11\quad 17 \]
Finn \os
\abch{
	\item typetallet \item medianen \item gjennomsnittet 
}
\newpage
\op{opgrikest}
Du ønsker å finne ut hva nordmenn flest har i formue\footnote{Enkelt sagt er formue summen av penger du har i banken, verdier av hus, bil etc., fratrekt gjeld o.l.}, og bestemmer deg for å finne ut av dette ved å spørre fem tilfeldige personer du møter i gata. De fire første svarene (i kr) er disse:
\[ 3,2\text{ millioner}\qquad 2,9\text{ millioner}\qquad 1,8\text{ millioner}\qquad 4,2\text{ millioner}  \]
Den siste personen du tilfeldigvis møter er personen i Norge med høyest formue. Hens svar er dette\footnote{Tallet er basert på personen i Norge med høyest formue ut ifra ligningstallene for 2019.}:
\[ 20\,915,3\text{ millioner} \]
\abc{
\item Finn medianen i datasettet.
\item Finn gjennomsnittet i datasettet.
\item Er det medianen eller gjennomsnittet som trolig best representerer hva nordmenn flest har i formue?
}

\op{opgfrkvtb1}
Lag en frekvenstabell for datasettet under. (La tittelen til venstre kolonne være ''frukt''.)
\begin{center}
	banan\quad eple \quad eple\quad eple \quad pære \quad banan \quad eple \quad pære \quad appelsin \quad eple \quad pære \quad pære
\end{center}

\op{opgfrkvtb2}
Lag en frekvenstabell for datasettet fra \refops{opgstatsentrpar2}. (La tittelen til venstre kolonne være ''tall''.)

\op{opgstatsoyl} \vs
\abc{
\item Lag et søylediagram for datasettet fra \refops{opgfrkvtb1}.
\item Lag et søylediagram for datasettet fra \refops{opgfrkvtb2}.
}

\op{opgstatsoylrgn} \vs
\abc{
	\item Lag et søylediagram for datasettet fra \refops{opgfrkvtb1}.
	\item Lag et søylediagram for datasettet fra \refops{opgfrkvtb2}.
}
\newpage

\op{opgstathvorfor}
Av de fire undersøkelsene på side \pageref{undersok}, hvorfor har vi
\abc{
	\item vist frekvenstabell bare for undersøkelse 2?
	\item vist søylediagram bare for undersøkelse  2 og 3?
	\item vist sektordiagram bare for undersøkelse 2 og ?
	\item vist linjediagram bare for undersøkelse 4?
}

\op{opgstat1}
Hvis datasettet har partalls antall svar kan man også finne medianen slik: \regv
\st{
\begin{enumerate}
	\item Finn de to tallene i midten.
	\item Finn differansen mellom tallene, og del denne med 2.
	\item Legg resultatet fra punkt 2 til det laveste av de to tallene i midten.
\end{enumerate}
} 
Prøv metoden på datasettet fra \refops{opgstatsentrpar1}.


\begin{comment}
	\op{opgstatsnittmobil}
	Tenk deg at noen i 2014 kom med følgende opplysning:
	\begin{center}
		''Mellom 2006 og 2014 ble det i gjennomsnitt solgt flere mobiltelefoner uten smartfunksjon enn med.''
	\end{center}
	Hvorfor ville denne opplysningen være lite til hjelp hvis man skulle spå framtidige salg av mobiltelefoner med og uten smartfunksjoner?
\end{comment}

\op{opgstathvorfor2}
Av de fire undersøkelsene på side \pageref{undersok}, hvorfor har vi ikke funnet sentral- og spredningsmål for undersøkelse 3?\vsk

\newpage
\eksop{GV26}{gv26d1opg1}
Diagrammet nedenfor viser antall arbeidsledige i Norge i aldersgruppen 15–24 år i perioden
2018–2024.
\abc{
\item Omtrent hvor mange menn var arbeidsledige i 2021?
\item Omtrent hvor mange færre arbeidsledige kvinner var det i 2022 sammenlignet med i 2021?
}
\begin{figure}
	\centering
	\includegraphics[scale=0.3]{\figp{gv26d1opg1}}
\end{figure}

\newpage
\grubeksoprecc{gv24d2opg8}{GV24D2}
En klasse fikk spørsmålet ''Hvor mange minutter bruker du i \\ gjennomsnitt på en dusj?''. Svarene ble fordelt på to datasett, ett for gutter og ett for jenter:
\begin{center}
	\small
	\begin{tabular}{l| c | c | c | c | c | c | c | c | c | c | c | c | c | c | c |}
		
		Gutter & 5 & 8 & 4 & 5 & 6 & 16 & 15 & 9 & 12 & 18 & 10 & 18 & 25 & 10 \\ \hline
		Jenter & 6 & 30 & 15 & 8 & 17 & 12 & 5 & 40 & 10 & 20 & 15 & 8 & 15 & 10 & 5
	\end{tabular}
\end{center}
\abc{
\item Finn typetallet (hvis det er et), medianen og gjennomsnittet for hver at datasettene.
\item Finn variasjonsbredden, variansen og standardavviket for hver av datasettene.
\item Kommenter hva målene du fant i oppgave a) og b) forteller oss om datasettene.
}

\grubop{opgstatfemsett}
Finn et datasett som oppfyller disse kriteriene:
\begin{itemize}
\item Det består av fem positive heltall
\item Gjennomsnittet er 4
\item Medianen er 3
\item Typetallet er 3
\end{itemize}

\grubeksoprecc{gv26d1opg10}{GV26D1}
Finn et datasett med fem heltall som oppfyller disse kriteriene:
\begin{itemize}
	\item Det består av fem positive heltall
	\item Gjennomsnittet er 5
	\item Typetallet er 6
	\item Medianen er 6
	\item Variasjonsbredden er 5
\end{itemize}


\begin{comment}
	\grubop{opgstatintrv}
	Ta utgangspunkt i tabellen på side \pageref{sportstable}.
	\abc{
		\item Lag et regneark med følgende overskrifter:
		\begin{center}
			\begin{tabular}{c|c|c|c|c}
				a & b & f & F & k
			\end{tabular}
		\end{center}
	}
\end{comment}

\grubop{gv23d1opg4} 
(GV23D1) \os
Tabellen nedenfor viser hastigheter målt i en
fartskontroll. Alle hastighetene er mål i km/h.
\begin{center}
	\begin{tabular}{|c|c|c|c|c|c|c|c|c|c|}
		62 & 20 & 62 & 18 & 55 & 62 & 65 & 54 & 62 & 60
	\end{tabular}
\end{center}
\abc{
\item Avgjør gjennomsnitt, median og typetall.
\item Begrunn hvilket av sentralmålene du ville valgt for å beskrive bilenes hastighet.
}

\eksop{G24}{gv24d1opg3}
Diagrammet under viser svarene avgitt av norske ungdommer i 2020 når de ble spurt om hvor mye tid de i gjennomsnitt brukte daglig på forskjellige medier. Tallene oppgir prosentandelen av svarene.
\begin{figure}
	\centering
	\includegraphics[scale=0.75]{\figp{gv24d1opg3}}
\end{figure}
Vurder om påstandene er sanne eller ikke sanne.
\begin{center}
	\begin{tabular}{| l | c | c |}
		\hline
		\textbf{Påstand} & \textbf{Sann} & \textbf{Ikke sann} \\ \hline	
	Ungdommer brukte mest tid på å lese bøker & & 	\\ \hline 
	Det var flere som spilte dataspill/TV-spill & & \\
	enn som så på TV. & & \\ \hline
	Nesten tre firedeler av ungdommene& & \\
	spilte på telefon/nettbrett & & \\ \hline
	Omtrent 40\,\% av ungdommene brukte & & \\
	mindre enn én time daglig på sosiale medier & & \\ \hline
	\end{tabular}
\end{center}

\eksop{G25}{gv25d1opg1}
Diagrammet nedenfor er fra en UngData-undersøkelse, og viser spisevaner til elever på
ungdomstrinnet. Tallene oppgir prosentandelen av svarene.
\fig{gv25d1opg1}
Vurder om påstandene er sanne eller usanne:
\begin{itemize}
	\item[(i)] Nesten $\dfrac{2}{3}$ av elevene spiser lunsj på skolen 5 dager i uka.
	\item[(ii)] Det er litt mer enn tre ganger så mange elever som spiser lunsj på
	skolen 5 dager i uka, enn de som spiser grønnsaker, frukt og bær på
	skolen 5 dager i uka.
	\item[(iii)] 60\,\% av ungdommene spiser grønnsaker, frukt eller bær 3 dager eller
	mer i uka.
	\item[(iv)] $\frac{1}{4}$ av elevene spiser lunsj på skolen 1-4 dager i uka.
\end{itemize}

\newpage
\grubeksoprecc{gv24d2opg1}{GV24D2}
Ida og Juan undersøker hva en ost, en pakke skinke og et brød koster til sammen i fire forskjellige matbutikker. De lager hvert sitt diagram som viser hva de forskjellige matvarene koster i de fire matbutikkene. 
\abc{
\item Avgjør hvilket diagram som best beskriver differansen mellom prisene. Grunngi svaret ditt.
\item Avgjør hvilket diagram som best kan brukes til å underbygge påstanden om at ''prisene i butikkene er relativt like''.
}
\fig{gb24d2opg1a}
\fig{gb24d2opg1b}

\newpage
\grubeksop{gv22d2opg5}{GV22D2}
Et politisk parti i Norge la frem grafen under for å vise\\ utviklingen av klimagassutslipp i Norge i perioden 2012-2018. \\ I denne perioden satt partiet i regjering fra og med 2013.
\abc{
\item Grafen fikk kritikk for å være misvisende. Gi minst ett\\ argument for at grafen er misvisende.
\item Lag et linjediagram hvor $y$-verdien til hvert punkt er\\ relativ til utslippet i 2012. (For eksempel: For 2014 vil den nye $y$-verdien være $\frac{\text{utslipp i 2014}}{\text{utslipp i 2012}}$). Gi minst ett argument for at dette linjediagrammet gir en bedre beskrivelse av utviklingen av klimagassutslippene.
}
\fig{gv22d2opg5}
\newpage
\grubop{gv23d1opg5}
(GV23D1) \os
Tabellen nedenfor viser hvor mange elever som bruker skoleskyss fordelt på fylke. \vs
\fig{gv23d1opg5}
Vurder om påstandene nedenfor er sanne eller usanne:
\begin{itemize}
	\item Flere enn tre ganger så mange elever bruker
	skoleskyss i \\ Innlandet som i Oslo.
	\item Gjennomsnittlig er det 5000 elever som bruker
	skoleskyss per fylke.
	\item I mer enn halvparten av fylkene er det under 25 %
	som bruker skoleskyss.
	\item Viken har den største prosentandelen av elever som bruker skoleskyss.
\end{itemize}

\newpage
\grubeksoprecc{gv25d2opg5}{GV25}
En skole samlet inn data over hvor mange timer hver elev brukte på
mobilen i løpet av en gjennomsnittlig dag. Tabellen viser svarene som ble gitt.
\begin{center}
	\begin{tabular}{l| c c c c c c c c c c c c c c c}
		Jenter &2 & 5 & 1 & 3 & 3 & 6 & 2 & 3 & 5 & 5 & 2 & 3 & 3 & 4 \\ \hline
		Gutter & 3 & 2 & 1 & 1 & 2 & 2 & 7 & 3 & 3 & 5 & 2 & 1 & 6 & 4 & 2
	\end{tabular}
\end{center}
\abc{
\item For jenter og gutter adskilt, og for klassen som helhet, finn typetallet, medianen, gjennomsnittet, variasjonsbredden, variansen og standardavviket.
\item Avgjør om de forskjellige medianene og gjennomsnittene er forholdsvis like eller ulike, og gi en forklaring av hvorfor det er slik.
\item Avgjør om standardavviket til klassen er forholdsvis stort eller ikke, og gi en forklaring av hvorfor det er slik.
\item Presenter tabellen i et diagram. Grunngi valg av diagram.
}

\grubeksoprecc{gv26d2opg1}{GV26}
Klasse 10C undersøkte hvor mange timer elevene i klassen sov én natt. Tabellen under viser
resultatet av undersøkelsen. (Hele timer skilt med komma.).
\begin{center}
	\begin{tabular}{l|l}
		Gutter & 7, 4, 4, 10, 6, 9, 2, 10, 12, 8, 8, 4 \\ \hline
		Jenter & 7, 6, 11, 6, 10, 9, 8, 8, 8, 9, 9, 5
	\end{tabular}
\end{center}
\abc{
\item For gutter og jenter adskilt, finn typetallet, gjennomsnittet, medianen, variasjonsbredden, variansen og standardavviket.
}
Anbefalt søvnmengde for en ungdom er 8 til 10 timer per natt. Elevene ønsker å finne ut hvor
stor andel av klassen som sov
\begin{itemize}
	\item mindre enn 8 timer
	\item 8 til 10 timer
	\item mer enn 10 timer
\end{itemize}
\abcs{2}{
\item Presenter dataene fra tabellen i et passende diagram, og bruk diagrammet til å
vurdere om elevene i 10 C får nok søvn.
}

\newpage
\grubeksoprecc{1pv26d1opg4}{1PV26D1}
Diagrammene under viser antall elger felt i fire kommuner i årene 2009-2019. Med \textit{økning fra ett år til det neste} mener vi økningen mellom et utvalgt år og året etter. Et eksempel er økningen\\ mellom 2016 og 2017 i Strand kommune.
\abc{
\item Hvilken kommune har hatt størst økning fra ett år til det neste målt i antall elger?
\item Hvilken kommune har hatt størst økning fra ett år til det neste målt i prosentvis endring?
}

\begin{centering}
	\includegraphics[scale=0.49]{\figp{1pv26d1opg4a}}
	\includegraphics[scale=0.49]{\figp{1pv26d1opg4b}} \\ [10pt]
	
	\includegraphics[scale=0.49]{\figp{1pv26d1opg4c}}
	\includegraphics[scale=0.49]{\figp{1pv26d1opg4d}} \\
\end{centering}
\end{document}

